\chapter{NEST: Implementation Details}
\label{chap:implementation-details}

Following the structural analysis and architectural designs established in Chapter 3, the implementation phase translates these logical specifications into executable, modular software components. This chapter provides a deep technical review of the NEST application's codebase, detailing the development stack, directory layouts, component design paradigms, reactive state flows, security mechanisms, and the implementations of the individual functional modules.

\section{Development Environment and Tooling}
To maximize development velocity, type safety, and runtime efficiency, the NEST application utilizes a modern, cross-platform software engineering stack:
The backend runtime is powered by Node.js (version 20.x), with npm serving as the package manager and scripts orchestrator. Strict TypeScript compilation rules are enforced across both frontend and backend projects (e.g., \texttt{"strict": true}, \texttt{"noImplicitAny": true}), catching type-coercion bugs and null reference pointer exceptions before the code is executed. On the frontend, the Angular CLI controls development compilation, serving, lazy-loaded chunk generation, and optimal Ahead-of-Time (AoT) production bundling. The database layer is managed by the Prisma Engine, which is instantiated as a development dependency to coordinate SQL schema synchronization, SQLite database seeding, and the generation of a strongly typed database client API. Finally, the SQLite Engine is utilized for relational persistence, running in-process to simplify local staging and deployment footprints while robustly supporting complex transaction models.

\section{Frontend Project Structure}
To sustain feature addition without architectural decay, the frontend application follows a strict Domain-Driven Design (DDD) layout. The directory organization within \texttt{frontend/src/app} enforces the separation of concerns:

\begin{enumerate}
    \item \textbf{\texttt{/core}:} Contains central application services that must only be instantiated once.
    \begin{itemize}
        \item \texttt{/constants}: Defines immutable configurations such as storage keys and API paths.
        \item \texttt{/guards}: Houses route-activation logic, preventing unauthenticated routing (\texttt{AuthGuard}) or block-pending routing restrictions (\texttt{PendingGuard}).
        \item \texttt{/interceptors}: Implements the \texttt{AuthInterceptor} which appends the user's JWT to all outgoing requests.
        \item \texttt{/services}: Declares foundational utilities such as \texttt{ThemeService}.
    \end{itemize}
    \item \textbf{\texttt{/store}:} Contains the centralized NgRx state architecture. Feature slices are isolated in individual subfolders (e.g., \texttt{auth}, \texttt{feed}, \texttt{shed}, \texttt{events}, \texttt{parking}) with their actions, reducers, effects, selectors, and facade definitions.
    \item \textbf{\texttt{/shared}:} Reusable assets shared across modules.
    \begin{itemize}
        \item \texttt{/components}: Standard presentational elements like spinners, modals, and headers.
        \item \texttt{/styles}: SCSS design tokens defining colors, layout, and grid dimensions.
    \end{itemize}
    \item \textbf{\texttt{/layout}:} Persistent framework structural shells, coordinating the sidebar, header, and toast container components.
    \item \textbf{\texttt{/features}:} Lazy-loaded feature modules corresponding to URL paths, containing routing definitions, visual page components, and local sub-components.
\end{enumerate}

\section{Angular Signal-Based Component Architecture}
NEST leverages the reactive state paradigm introduced in modern Angular releases, using Signal APIs rather than legacy dirty-checking engines. This change ensures precise DOM updates, as the framework directly identifies which node in the template depends on a modified Signal value, avoiding full-page re-renders \cite{angularSignalsRfc}.

To illustrate this architecture, Listing \ref{lst:page-header} shows the TypeScript implementation of the \texttt{PageHeaderComponent} (\texttt{frontend/src/app/shared/components/page-header/page-header.component.ts}), which utilizes standalone component declaration, decoupled stylesheet and template references, and modern Signal inputs.

\begin{lstlisting}[caption={Page Header Component with Separated TS Controller}, label={lst:page-header}]
import { Component, input } from '@angular/core';

@Component({
  selector: 'app-page-header',
  standalone: true,
  templateUrl: './page-header.component.html',
  styleUrl: './page-header.component.scss'
})
export class PageHeaderComponent {
  readonly title = input.required<string>();
  readonly subtitle = input<string>('');
}
\end{lstlisting}

By utilizing \texttt{input.required<string>()} and \texttt{input<string>('')}, properties are exposed as read-only Signals. Component templates (separated into \texttt{page-header.component.html}) read these values using function invocation syntax (e.g., \texttt{title()}), establishing a dependency link that automatically updates the specific template node when the underlying input changes.

\section{Reactive State Management via NgRx and Facades}
To manage complex, asynchronous, multi-user application state without cluttering components, NEST implements a global Redux architecture through NgRx \cite{ngrxDocs}. The client-side state is centralized in a single store, updated through a unidirectional data flow: components trigger actions, reducers perform synchronous state updates, effects handle asynchronous HTTP network requests, and selectors provide components with sliced state views.

To decouple components from store implementation details, the architecture applies the **Facade Pattern** \cite{gamma1994design}. In NEST, the Facade acts as the sole gateway for a feature's state. Rather than having visual components understand the internal structure of the NgRx store, dispatch specific raw actions, or select complex state slices, the components simply call methods like \texttt{authFacade.login()} or read signals like \texttt{authFacade.currentUser()}. The Facade internally translates these calls into the appropriate NgRx \texttt{Store.dispatch} operations. Listing \ref{lst:auth-facade} displays the implementation of the \texttt{AuthFacade} (\texttt{frontend/src/app/store/auth/auth.facade.ts}), illustrating how it exposes read-only reactive properties as Signals using \texttt{selectSignal()}.

\begin{lstlisting}[caption={Auth Facade isolating NgRx Store mechanics from Components}, label={lst:auth-facade}]
import { Injectable, inject } from '@angular/core';
import { Store } from '@ngrx/store';
import { AuthActions } from './auth.actions';
import {
  selectCurrentUser,
  selectToken,
  selectPermissions,
  selectAuthIsLoading,
  selectAuthError,
  selectIsAuthenticated,
  selectIsVerified,
  selectIsAdmin,
  selectUserFullName,
} from './auth.selectors';
import { User } from '../../core/models/user.model';
import { TOKEN_STORAGE_KEY, USER_STORAGE_KEY } from '../../core/constants/ui';

@Injectable({ providedIn: 'root' })
export class AuthFacade {
  private store = inject(Store);

  currentUser = this.store.selectSignal(selectCurrentUser);
  token = this.store.selectSignal(selectToken);
  permissions = this.store.selectSignal(selectPermissions);
  isLoading = this.store.selectSignal(selectAuthIsLoading);
  error = this.store.selectSignal(selectAuthError);
  isAuthenticated = this.store.selectSignal(selectIsAuthenticated);
  isVerified = this.store.selectSignal(selectIsVerified);
  isAdmin = this.store.selectSignal(selectIsAdmin);
  userFullName = this.store.selectSignal(selectUserFullName);

  login(email: string, password: string): void {
    this.store.dispatch(AuthActions.login({ email, password }));
  }

  register(email: string, password: string, firstName: string, lastName: string, apartmentNumber: string): void {
    this.store.dispatch(AuthActions.register({ email, password, firstName, lastName, apartmentNumber }));
  }

  joinBlock(blockCode: string): void {
    this.store.dispatch(AuthActions.joinBlock({ blockCode }));
  }

  logout(): void {
    this.store.dispatch(AuthActions.logout());
  }

  tryRestoreSession(): void {
    const token = localStorage.getItem(TOKEN_STORAGE_KEY);
    const userJson = localStorage.getItem(USER_STORAGE_KEY);

    if (token && userJson) {
      const user: User = JSON.parse(userJson);
      this.store.dispatch(AuthActions.restoreSession({ user, token }));
    }
  }

  updateUser(user: User): void {
    this.store.dispatch(AuthActions.updateUser({ user }));
    localStorage.setItem(USER_STORAGE_KEY, JSON.stringify(user));
  }
}
\end{lstlisting}

This facade pattern provides three significant engineering advantages:
\begin{enumerate}
    \item \textbf{Simplified Component Injection:} Components only need to inject \texttt{AuthFacade} rather than injecting \texttt{Store}, importing various actions, and injecting multiple selectors.
    \item \textbf{Signal-Based Interoperability:} Exposing properties as Signals allows components to use Angular's fine-grained reactivity directly inside their templates and controller \texttt{effect()} bindings.
    \item \textbf{Store Isolation:} If the underlying state library is replaced (e.g., migrating from NgRx to NgRx Signals Store), the visual components remain completely unchanged, preserving the codebase's maintainability.
\end{enumerate}

\section{Applied Design Patterns}

Several classical design patterns from the ``Gang of Four'' catalog \cite{gamma1994design} find direct application in the architecture of NEST, ensuring a maintainable and scalable codebase.

\textbf{Singleton.} Angular's dependency injection system guarantees that services declared with \texttt{providedIn: 'root'} are instantiated exactly once and shared across the entire application. This is utilized in NEST for the \texttt{ThemeService} and \texttt{AuthFacade}, where multiple instances would produce conflicting states (e.g., UI flickering between themes or desynced authentication states).

\textbf{Observer.} The NgRx store is fundamentally an implementation of the Observer pattern. Components \textit{subscribe} to slices of the state, and the store \textit{notifies} them whenever the relevant state changes. In NEST's reactive programming model, this is expressed through Angular Signals (e.g., \texttt{this.store.selectSignal}) and RxJS Observables, ensuring that the Feed and Events modules update in real-time when new data arrives.

\textbf{Strategy.} The Express middleware pipeline on the backend exemplifies the Strategy pattern. Each middleware function (\texttt{authenticateToken}, \texttt{verifyRole}, \texttt{rateLimiter}) encapsulates a specific policy. They are composed dynamically depending on the route's security requirements. For example, the Parking creation route uses both the authentication strategy and the resident-role strategy.

\textbf{Dependency Injection.} Angular's DI container is perhaps the most pervasive pattern in the framework. Services, guards, interceptors, and resolvers are all instantiated and injected by the framework based on their provider configuration. In the NEST frontend, this allows the \texttt{AuthInterceptor} to seamlessly inject the \texttt{AuthFacade} to retrieve the current JWT token for outgoing requests without hardcoding dependencies.

\section{Core Feature Modules Implementation}
NEST's user experience is structured across 8 feature modules, each designed to address a specific aspect of localized community interaction.

\subsection{Authentication and Dynamic Block Onboarding}
The registration and login flow coordinates physical building associations through a multi-stage onboarding system:
\textbf{Algorithm: Multi-Stage Registration and Join Flow}
\begin{enumerate}[label=\textbf{Step \arabic*:}, leftmargin=*]
    \item \textbf{Initialize Registration:} The guest submits account details (name, email, password, apartment). The backend creates a User record with a \texttt{PENDING} status, generates a JWT, and authenticates the session.
    \item \textbf{Intercept Routing:} The frontend \texttt{PendingGuard} intercepts all navigation attempts. Identifying that the user's \texttt{blockId} is null, the guard redirects the user to the Block Selection interface.
    \item \textbf{Submit Join Request:} The user inputs their physical building's unique alphanumeric code. The system validates the code against the database and creates a \texttt{JoinRequest} entity linked to the user and the block.
    \item \textbf{Admin Authorization:} The block administrator logs into the Admin Panel, reviews the pending \texttt{JoinRequest} list, verifies the resident's physical identity, and triggers the approval endpoint.
    \item \textbf{Activate Session:} The backend updates the User's \texttt{isVerified} flag to true and assigns the \texttt{blockId}. The \texttt{PendingGuard} clears the routing restriction, granting the user full access to community modules.
\end{enumerate}

\subsection{The Community Feed: Image Upload and Pagination}
The Community Feed serves as the digital billboard for neighborhood announcements:
\begin{itemize}
    \item \textbf{Data Fetching:} Utilizes Angular Route Resolvers to pre-fetch feed data before rendering components, preventing page layout shifts.
    \item \textbf{Image Storage:} The backend uses Multer middleware to process image attachments, storing files in static upload directories and saving path URLs in the database's \texttt{Post.imageUrl} column.
    \item \textbf{Reactivity:} Leverages NgRx Effects to append newly published posts directly to the frontend store, refreshing the UI without requiring full manual page refreshes.
\end{itemize}

\subsection{The Shared Shed: Peer-to-Peer Booking State Machine}
The Shared Shed implements peer-to-peer resource sharing. Items listed in the shed follow a strict state-machine lifecycle managed through transactional database writes:

\begin{figure}[h]
\centering
\begin{tikzpicture}[node distance=2.4cm,
    state/.style={draw, rectangle, minimum height=0.8cm, minimum width=2cm,
                  align=center, font=\small},
    lbl/.style={font=\tiny, fill=white, inner sep=1pt},
    >=Stealth]
    \node[state] (avail) {AVAILABLE};
    \node[state, right=of avail] (pending) {PENDING};
    \node[state, right=of pending] (approved) {APPROVED};
    \node[state, right=of approved] (returned) {RETURNED};
    \draw[->] (avail) -- node[lbl,above] {Request} (pending);
    \draw[->] (pending) -- node[lbl,above] {Approve} (approved);
    \draw[->] (approved) -- node[lbl,above] {Return} (returned);
    \draw[->] (returned.north) -- ++(0,0.6) -| node[lbl, pos=0.25, above] {Make Available Again} (avail.north);
\end{tikzpicture}
\caption{State transition lifecycle of a Shared Shed resource reservation.}
\label{fig:resource-state-machine}
\end{figure}

To prevent race conditions where multiple users try to reserve the same resource for overlapping periods, the Express service utilizes Prisma's relational transaction API:
\begin{lstlisting}
const activeOverlaps = await prisma.resourceReservation.count({
  where: {
    resourceId: req.body.resourceId,
    status: 'APPROVED',
    OR: [
      { startTime: { lte: req.body.endTime }, endTime: { gte: req.body.startTime } }
    ]
  }
});
if (activeOverlaps > 0) throw new Error('Resource already reserved during this interval');
\end{lstlisting}

\subsection{Events Module: Participant Limits and RSVP Control}
The Events module allows residents to plan physical activities. A critical technical detail is the dynamic enforcement of attendee limits:
\begin{itemize}
    \item \textbf{Data Structure:} Tracked using the \texttt{EventAttendee} join table, which maps users to events.
    \item \textbf{Verification:} When a user clicks "RSVP Join", the backend service counts active participants. If the current participant count is greater than or equal to the event's \texttt{maxParticipants} capacity limit, the API rejects the request, returning a \texttt{400 Bad Request} validation error.
    \item \textbf{Reactive Feedback:} The UI updates dynamically, rendering disabled buttons and "Event Full" indicators in real-time as state changes are propagated.
\end{itemize}

\subsection{Parking Sharing Module: Conflict Resolution}
The Parking module manages temporary parking slot allocations. Spot owners list their spaces for specific intervals, and other residents apply to rent them:
\begin{enumerate}
    \item \textbf{Conflict Prevention:} To maintain order, the backend implements auto-conflict resolution. 
    \item \textbf{Application Approval:} When an owner approves a specific application (\texttt{status\allowbreak\ =\allowbreak\ "APPROVED"}), the service triggers a database transaction that automatically marks all other applications for that announcement as \texttt{DENIED}.
    \item \textbf{Payload Response:} This cascading update ensures slots cannot be double-booked, sending notifications to all affected applicants and updating their UI states.
\end{enumerate}

\section{CSS Design Token System and Theming Engine}
To provide a cohesive visual experience, the application avoids ad-hoc styles in favor of a centralized design token system. All styling dimensions, layouts, and colors are defined using CSS Custom Properties (Variables) inside the \path{frontend/src/app/shared/styles/_colors.scss} stylesheet \cite{mdnCssCustomProperties}. 

The architecture supports light and dark themes using a root dataset selector. Theme variables are mapped onto semantic CSS tokens under the \texttt{:root} element, and overridden under the \texttt{[data-theme='dark']} attribute selector. Listing \ref{lst:colors-excerpt} shows an excerpt of the design system variables:

\begin{lstlisting}[caption={Color Design Token overrides for Light and Dark themes}, label={lst:colors-excerpt}]
:root {
  --color-emerald-50: #ecfdf5;
  --color-emerald-500: #10b981;
  --color-emerald-900: #064e3b;
  
  --color-bg: #ecfdf5;
  --color-surface: #ffffff;
  --color-text-primary: #111827;
  --color-border: #d1fae5;
  --color-primary: #059669;
}

[data-theme='dark'] {
  --color-bg: #030712;
  --color-surface: #111827;
  --color-text-primary: #f9fafb;
  --color-border: #1f2937;
  --color-primary: #10b981;
}
\end{lstlisting}

Theme switching is managed reactively via the \texttt{ThemeService} (\path{frontend/src/app/shared/services/theme.service.ts}). It uses an Angular Signal to track the current theme state and an Angular \texttt{effect()} to write the active class attribute to the document body, persisting the user's choice in browser storage:

\begin{lstlisting}
@Injectable({ providedIn: 'root' })
export class ThemeService {
  private activeTheme = signal<string>('light');

  constructor() {
    const saved = localStorage.getItem('theme');
    if (saved) this.activeTheme.set(saved);
    
    effect(() => {
      const current = this.activeTheme();
      document.body.setAttribute('data-theme', current);
      localStorage.setItem('theme', current);
    });
  }

  toggleTheme(): void {
    this.activeTheme.update(t => t === 'light' ? 'dark' : 'light');
  }
}
\end{lstlisting}

By utilizing this token-based architecture, components do not define hardcoded hex colors. Instead, they reference semantic tokens (e.g., \path{color: var(--color-text-primary)}). When the theme attribute changes, the browser updates colors across the DOM tree automatically, providing a smooth user experience.

\section{Backend Security Implementation}
To safeguard resident data and prevent unauthorized API access, NEST implements a multi-layered security pipeline. Security policies are enforced before requests reach controllers through a chain of modular Express middleware:

\begin{lstlisting}[caption={Express Security Middleware Chain}, label={lst:security-middleware}]
import { Request, Response, NextFunction } from 'express';
import { verifyToken } from '../utils/jwtUtils';
import { ROLES } from '../config/constants';

export interface AuthenticatedRequest extends Request {
  user?: any;
}

export const requireAuthentication = (req: AuthenticatedRequest, res: Response, next: NextFunction): void => {
  const authorizationHeader = req.headers.authorization;
  if (!authorizationHeader || !authorizationHeader.startsWith('Bearer ')) {
    res.status(401).json({ error: 'Unauthorized credentials' });
    return;
  }
  const token = authorizationHeader.split(' ')[1];
  const decoded = verifyToken(token);
  if (!decoded) {
    res.status(401).json({ error: 'Expired or invalid credentials' });
    return;
  }
  req.user = decoded;
  next();
};

export const requireVerified = (req: AuthenticatedRequest, res: Response, next: NextFunction): void => {
  if (!req.user) {
    res.status(401).json({ error: 'Authentication required' });
    return;
  }
  if (!req.user.isVerified) {
    res.status(403).json({ error: 'Account not yet verified by block admin' });
    return;
  }
  next();
};

export const requireAdminRole = (req: AuthenticatedRequest, res: Response, next: NextFunction): void => {
  if (!req.user) {
    res.status(401).json({ error: 'Authentication required' });
    return;
  }
  if (req.user.role !== ROLES.ADMIN) {
    res.status(403).json({ error: 'Insufficient administrative privileges' });
    return;
  }
  next();
};
\end{lstlisting}

In the routing definitions, these middleware checks are chained sequentially to enforce different security boundaries:
\begin{itemize}
    \item \textbf{General Security:} \texttt{router.get('/feed',} \allowbreak \texttt{requireAuthentication,} \allowbreak \texttt{requireVerified,} \allowbreak \texttt{feedController.getFeed)} protects standard community data.
    \item \textbf{Administrative Restrictions:} \texttt{router.put('/admin/approve/:id',} \allowbreak \texttt{requireAuthentication,} \allowbreak \texttt{requireAdminRole,} \allowbreak \texttt{adminController.approveUser)} restricts administrative actions to authorized accounts.
\end{itemize}

This multi-layered approach ensures that unauthenticated clients are blocked at the perimeter (\texttt{requireAuthentication}), pending registrations are restricted from standard access (\texttt{requireVerified}), and administrative actions are isolated to verified administrators (\texttt{requireAdminRole}).